\documentclass[preview, border={10mm, 5mm, 10mm, 10mm}]{standalone}
\usepackage{tikz, tasks, geometry, xcolor}
\usetikzlibrary{arrows.meta, patterns, calc, intersections, quotes, angles}
\usepackage{amsmath, amssymb, amsfonts, mathtools}
\setlength{\columnsep}{10pt}
\setlength{\columnseprule}{0.4pt}
\usepackage[upright]{fourier}
\usepackage{enumitem}
\geometry{a4paper, margin=1in}
\usepackage{tzplot, pgfplots, kinematikz}
\usepackage{circuitikz}
\ctikzset{resistors/scale=0.75,capacitors/scale=0.75,inductors/scale=0.75}
\usepackage{chemfig}
\usepackage[version=4]{mhchem}
\usepackage[modules=all]{chemmacros}
\usepackage{tkz-euclide}
\usepackage{venndiagram}
\pgfplotsset{compat=1.18}
\everymath{\displaystyle}
\newcommand{\ans}{\textcolor{blue!20!red}{\textit{\quad Ans.}}}
% \renewcommand{\ans}{}  % Uncomment to hide answers
\newenvironment{solution}{\par\noindent\color{red!80!black}$\Rightarrow$\enspace\ignorespaces}{\par}
\newenvironment{alternatesolution}{\par\noindent\color{blue!80!black}$\Rrightarrow$\enspace\ignorespaces}{\par}
\newenvironment{hint}{\par\noindent\color{red!50!black}$\looparrowright$\enspace\ignorespaces}{\par}
\newenvironment{idea}{\par\noindent\color{violet!80!black}$\diamond$\enspace\ignorespaces}{\par}
\newenvironment{remark}{\par\noindent\color{teal!80!black}$\circ$\enspace\ignorespaces}{\par}

% --- Global TikZ style (design uniformity across all diagrams) ---
\tikzset{
    >=latex,
    thick,
    every node/.append style={font=\small},
}
\title{\textsc{}}
\pagestyle{empty}
\begin{document}
\begin{enumerate}[leftmargin=*]
\item[$\Omega.$] An electron (mass $9 \times 10^{-31}\,\mathrm{kg}$ and charge $1.6 \times 10^{-19}\,\mathrm{C}$) moving with speed $c/100$ ($c$ = speed of light) is injected into a magnetic field $\vec{B}$ of magnitude $9 \times 10^{-4}\,\mathrm{T}$ perpendicular to its direction of motion. We wish to apply an uniform electric field $\vec{E}$ together with the magnetic field so that the electron does not deflect from its path. Then (speed of light $c = 3 \times 10^8\,\mathrm{m s}^{-1}$)
\begin{tasks}(1)
    \task $\vec{E}$ is parallel to $\vec{B}$ and its magnitude is $27 \times 10^2\,\mathrm{V m}^{-1}$
    \task $\vec{E}$ is parallel to $\vec{B}$ and its magnitude is $27 \times 10^4\,\mathrm{V m}^{-1}$
    \task $\vec{E}$ is perpendicular to $\vec{B}$ and its magnitude is $27 \times 10^4\,\mathrm{V m}^{-1}$
    \task $\vec{E}$ is perpendicular to $\vec{B}$ and its magnitude is $27 \times 10^2\,\mathrm{V m}^{-1}$ 
\end{tasks}
\end{enumerate}
\hfill $\_\_$
\end{document}
